% =============================================================================
% P8 cost-per-decision accounting + LLM-as-sensor feature ablation -- iter 124
% =============================================================================
\subsection{Cost-per-decision accounting on the iter-120 V\_stat ablation
  (\textit{iter 124})}
\label{sec:p8-iter124-cost-accounting}

The iter-80 gradient-band rule and the iter-108/iter-120 cohort-asymmetry
analyses reported a single LLM-call price point (\$0.001/call,
COST\_LLM/COST\_XGB = 10$\times$). Real 2026 LLM API prices vary by
500$\times$ (heuristic rule at \$0.0001/call; small open LLM at
\$0.0006/call; mid-tier at \$0.005/call; frontier GPT-4-class at
\$0.030/call).  Iter-124 sweeps these five price tiers and reports
the per-$V$\_stat quartile {\bf cppr} (cost per positive RECALLED) and
{\bf cpd} (cost per DECISION).

\subsubsection{Falsifiable headline H1 -- grad-band is NOT cheaper at any
  realistic LLM tier}
On the iter-120 24full backbone, the cppr ratio
(cppr$_{\text{grad}}$ $/$ cppr$_{\text{xgb}}$) at the five price tiers:
\begin{itemize}
\item cheap\_heuristic (\$0.0001$/$call): ratio = 1.000 (tie)
\item small\_open    (\$0.0006$/$call): ratio = 1.041
\item iter120\_default (\$0.001$/$call): ratio = 1.075
\item mid\_tier      (\$0.005$/$call):  ratio = 1.412
\item frontier\_gpt4 (\$0.030$/$call):  ratio = 3.512
\end{itemize}
\textbf{Verdict (H1 REFUTED)} -- the gradient-band rule is strictly
cheaper than xgb-only ONLY at the trivial cost$_{\text{LLM}}$ = cost$_{\text{XGB}}$
price tier.  At all four realistic price tiers, grad-band has HIGHER cppr.
The honest framing: \emph{the iter-80 rule's value is NOT a cost saving --
it is a recall-augmentation signal that recovers xgb-only-missed positives
on the 0.84\% most uncertain top-K rows.}

\subsubsection{Falsifiable headline H2 -- worst-cost $V$\_stat quartile is
  STABLE across price tiers}
For each price tier, the $V$\_stat quartile with the highest
cppr$_{\text{grad}}$ $/$ cppr$_{\text{xgb}}$ ratio is $V$\_mean $Q_2$ at
four of five tiers (small\_open through frontier\_gpt4).
$V$\_mean $Q_2$ has the highest call density (37 LLM fires on 14
xgb-only-caught positives) and the lowest xgb-only recall (0.28).
The stable worst-cost cell is a structural property of the iter-80
rule, not a cost artefact.

\subsubsection{Falsifiable headline H3 -- LLM-as-sensor feature ablation}
Retrained XGBoost with four feature subsets:
\begin{itemize}
\item 24full (24 feats): xgb-only recall$@K{=}2\%$ = 53$/$144 = 0.368, n$_{\text{llm}}$ = 84
\item 20raw (20 feats): xgb-only recall$@K{=}2\%$ = 53$/$144 = 0.368, n$_{\text{llm}}$ = 81
\item 20raw$+$minmax (22 feats): xgb-only recall$@K{=}2\%$ = 50$/$144 = 0.347, n$_{\text{llm}}$ = 84
\item 20raw$+$stat (22 feats): xgb-only recall$@K{=}2\%$ = 50$/$144 = 0.347, n$_{\text{llm}}$ = 85
\end{itemize}
\textbf{Verdict (H3 PASS)} -- dropping the four aggregate features
($V$\_mean, $V$\_std, $V$\_max, $V$\_min) changes the gradient-band firing
pattern by $-3$ calls (24full$\to$20raw, 96\% preserved). Pairwise
agreement between 24full and 20raw firing masks is $\geq$0.99.
The rule captures a property of the score stream
(consecutive-gradient sparsity in the top-K), not a property of the
model input space.

\subsubsection{Falsifiable headline H4 -- sweet-spot price per quartile}
Closed-form sweet-spot price per $V$\_stat quartile:
$\text{price}_{\text{sweet}} = \text{cost}_{\text{xgb}} \cdot
\text{xgb\_caught}_{q} / n_{\text{llm},q}$.  Range across 16 quartiles:
$\$0.000017$ ($V$\_std $Q_0$) to $\$0.000120$ ($V$\_std $Q_3$).
All sweet-spot prices are below cost$_{\text{xgb}} = \$0.0001$ baseline,
confirming that the rule is cost-rational only at the trivial
cost$_{\text{LLM}} =$ cost$_{\text{XGB}}$ price tier.

\subsubsection{Operational recommendation}
The iter-80 gradient-band rule's role is to \textbf{recover xgb-only-missed
positives on the 0.84\% most uncertain top-K rows}.  The cost of doing so
is real (1.04$\times$ to 3.51$\times$ higher than xgb-only on cppr,
depending on LLM tier), and the cost is the price of catching the residual.
\textbf{Do not frame the rule as a cost optimizer} -- it is a recall
augmentation.  The honest framing for the paper is the
\emph{cost-rationing} view: the rule is cost-rational only at the
cost$_{\text{LLM}} \leq \$0.0001$ tier; at \$0.005 (mid-tier), the rule
is 41$\times$ over its sweet-spot price.